\documentclass[12pt,a4paper]{article}
\usepackage[utf8]{inputenc}
\usepackage{amsmath}
\usepackage{amssymb}
\usepackage{algorithm}
\usepackage{algpseudocode}
\usepackage{listings}
\usepackage{xcolor}
\usepackage{geometry}
\usepackage{graphicx}
\usepackage{tcolorbox}
\usepackage{enumitem}
\usepackage{longtable}
\usepackage{array}

\geometry{margin=0.5in}

% Code listing style
\lstset{
    basicstyle=\ttfamily\small,
    breaklines=true,
    frame=single,
    backgroundcolor=\color{gray!10},
    keywordstyle=\color{blue},
    commentstyle=\color{green!50!black},
    stringstyle=\color{red},
    numbers=left,
    numberstyle=\tiny\color{gray},
    stepnumber=1,
    numbersep=5pt,
    showstringspaces=false
}

\title{\textbf{Two-Stage Stochastic Nurse Scheduling with CVaR Risk Control}\\
\large Mathematical Model and Algorithms}
\author{Nurse Scheduling System (NSS)\\
Based on He, Chaussalet \& Qu (2019)}
\date{\today}

\begin{document}

\maketitle

\tableofcontents
\newpage

\section{Introduction}

This document presents the complete mathematical formulation of the integrated nurse staffing and scheduling problem under patient demand uncertainty. The model is formulated as a two-stage stochastic integer program with Conditional Value-at-Risk (CVaR) constraints for understaffing risk control.

\subsection{Problem Overview}

\begin{itemize}
    \item \textbf{Stage 1 (Here-and-Now):} Decide baseline nurse schedule before demand is known
    \item \textbf{Stage 2 (Recourse):} Adjust schedule after demand realization by adding/cancelling shifts
    \item \textbf{Risk Control:} Use CVaR to limit worst-case understaffing scenarios
\end{itemize}

\section{Notation}

\subsection{Sets and Indices}

\begin{tcolorbox}[title=Sets]
\begin{align*}
I &= \text{Set of nurses} \quad (i \in I) \\
J &= \text{Set of days in planning period} \quad (j \in J) \\
K &= \text{Set of shift types: } \{E, D, L, N\} \quad (k \in K) \\
&\quad \text{E = Early, D = Day, L = Late, N = Night} \\
\Omega &= \text{Set of demand scenarios} \quad (\omega \in \Omega) \\
W &= \text{Set of weeks in planning period} \quad (w \in W) \\
S &= \text{Set of PWL segments for fatigue approximation} \quad (s \in S) \\
K' &= \text{Set of unwanted shift patterns} \\
&\quad K' = \{(D,E), (L,E), (L,D), (E,N)\}
\end{align*}
\end{tcolorbox}

\subsection{Parameters}

\begin{tcolorbox}[title=Stage 1 Parameters]
\begin{align*}
c_1 &= \text{Regular shift cost per shift} \\
c_2 &= \text{Overtime shift cost per shift} \\
c_3 &= \text{Penalty for stand-alone shift violation} \\
c_4 &= \text{Penalty for unwanted shift pattern violation} \\
c_{safety} &= \text{Penalty per unit of fatigue (Patient Safety Weight)} \\
\lambda &= \text{Fatigue accumulation rate (from Jaber et al. 2013)} \\
n_1 &= \text{Maximum total shifts per nurse} \\
n_2 &= \text{Maximum night shifts per nurse} \\
n_3 &= \text{Minimum regular shifts per nurse} \\
n_4 &= \text{Minimum complete weekends off per nurse} \\
R_{jk} &= \text{Baseline demand for shift } k \text{ on day } j
\end{align*}
\end{tcolorbox}

\begin{tcolorbox}[title=Stage 2 Parameters]
\begin{align*}
R_{jk}^\omega &= \text{Realized demand for shift } k \text{ on day } j \text{ in scenario } \omega \\
p^\omega &= \text{Probability of scenario } \omega \quad \left(\sum_{\omega \in \Omega} p^\omega = 1\right) \\
q^+ &= \text{Cost of adding one emergency shift} \\
q^- &= \text{Cost of cancelling one shift} \\
\alpha^{max} &= \text{Maximum emergency staff per shift (Constraint 17)} \\
\beta^{max} &= \text{Maximum cancellations per shift (Constraint 18)}
\end{align*}
\end{tcolorbox}

\begin{tcolorbox}[title=CVaR Parameters]
\begin{align*}
\sigma &= \text{Confidence level (e.g., 0.95 for 95\%)} \\
\mu &= \text{Maximum acceptable shortage threshold} \\
F_{max} &= \text{Maximum allowed fatigue threshold (Project Constraint)}
\end{align*}
\end{tcolorbox}

\subsection{Decision Variables}

\begin{tcolorbox}[title=Stage 1 Variables (First-Stage)]
\begin{align*}
sr_{ijk} &\in \{0,1\} \quad \text{= 1 if nurse } i \text{ works regular shift } k \text{ on day } j \\
so_{ijk} &\in \{0,1\} \quad \text{= 1 if nurse } i \text{ works overtime shift } k \text{ on day } j \\
SR_i &\in \{0,1\} \quad \text{= 1 if nurse } i \text{ works any regular shifts} \\
SO_i &\in \{0,1\} \quad \text{= 1 if nurse } i \text{ works any overtime shifts} \\
F_{ij} &\in [0, F_{max}] \quad \text{Cumulative fatigue for nurse } i \text{ on day } j \\
T_{ij} &\in \mathbb{R}_+ \quad \text{Cumulative work hours for nurse } i \text{ up to day } j \\
w_{ijs} &\in [0,1] \quad \text{Weight for PWL breakpoint } s \quad (\text{SOS2 Logic}) \\
y_{ijs} &\in \{0,1\} \quad \text{Segment selection indicator for fatigue PWL} \\
dev1_{ij} &\in \mathbb{Z}_+ \quad \text{Deviation for stand-alone shift on day } j \\
dev2_{ijk} &\in \mathbb{Z}_+ \quad \text{Deviation for unwanted pattern starting with shift } k
\end{align*}
\end{tcolorbox}

\begin{tcolorbox}[title=Stage 2 Variables (Recourse)]
\begin{align*}
\alpha_{jk}^\omega &\in \mathbb{Z}_+ \quad \text{Number of shifts added on day } j, \text{ shift } k, \text{ scenario } \omega \\
\beta_{jk}^\omega &\in \mathbb{Z}_+ \quad \text{Number of shifts cancelled on day } j, \text{ shift } k, \text{ scenario } \omega
\end{align*}
\end{tcolorbox}

\begin{tcolorbox}[title=CVaR Variables]
\begin{align*}
\xi &\in \mathbb{R} \quad \text{Value-at-Risk (VaR) threshold} \\
z^\omega &\in \mathbb{R}_+ \quad \text{Excess loss in scenario } \omega \text{ beyond VaR}
\end{align*}
\end{tcolorbox}

\newpage

\section{Objective Function}

The objective is to minimize the total expected cost, which consists of:

\begin{tcolorbox}[title=Complete Objective Function]
\begin{equation}
\begin{aligned}
\min \quad & \underbrace{c_1 \sum_{i \in I} \sum_{j \in J} \sum_{k \in K} sr_{ijk}}_{\text{Regular shift cost}} 
+ \underbrace{c_2 \sum_{i \in I} \sum_{j \in J} \sum_{k \in K} so_{ijk}}_{\text{Overtime shift cost}} \\
& + \underbrace{c_3 \sum_{i \in I} \sum_{j \in J} dev1_{ij}}_{\text{Stand-alone penalty}}
+ \underbrace{c_4 \sum_{i \in I} \sum_{j \in J} \sum_{k \in K} dev2_{ijk}}_{\text{Unwanted pattern penalty}} \\
& + \underbrace{c_{safety} \sum_{i \in I} \sum_{j \in J} F_{ij}}_{\text{Safety Risk (Project Enhancement)}} \\
& + \underbrace{\sum_{\omega \in \Omega} p^\omega \left( q^+ \sum_{j \in J} \sum_{k \in K} \alpha_{jk}^\omega + q^- \sum_{j \in J} \sum_{k \in K} \beta_{jk}^\omega \right)}_{\text{Expected recourse cost}}
\end{aligned}
\end{equation}
\end{tcolorbox}

\subsection{Python Implementation}

\begin{lstlisting}[language=Python, caption=Objective Function Implementation (model.py:L615)]
    # STAGE 1 COST: Regular wages + Overtime wages
    stage1_cost = (
        c1 * pulp.lpSum(sr[i][j][k] for i in I_nurses for j in J_days for k in K_shifts) +
        c2 * pulp.lpSum(so[i][j][k] for i in I_nurses for j in J_days for k in K_shifts)
    )
    
    # SOFT CONSTRAINT PENALTIES
    soft_penalty_cost = (
        c3 * pulp.lpSum(dev1[i][j] for i in I_nurses for j in J_days) +
        c4 * pulp.lpSum(dev2[i][j][k] for i in I_nurses for j in J_days for k in K_shifts)
    )

    # STAGE 2 COST: Expected recourse cost across all scenarios
    stage2_cost = pulp.lpSum(
        scenario_probability[w] * (
            q_plus * pulp.lpSum(alpha[j][k][w] for j in J_days for k in K_shifts) +
            q_minus * pulp.lpSum(beta[j][k][w] for j in J_days for k in K_shifts)
        )
        for w in W_scenarios
    )
    
    # PATIENT SAFETY COST: Fatigue-related costs
    if patient_safety_enabled:
        patient_safety_cost = patient_safety_weight * pulp.lpSum(F[i][j] for i in I_nurses for j in J_days)
    else:
        patient_safety_cost = 0
    
    # TOTAL OBJECTIVE
    prob += stage1_cost + soft_penalty_cost + stage2_cost + patient_safety_cost, "Total_Cost"
\end{lstlisting}

\newpage

\section{Constraints}

\subsection{Stage 1 Constraints (Baseline Schedule)}

\subsubsection{Constraint 1: One Shift Per Day}

\begin{tcolorbox}[title=Constraint 1: Maximum One Shift Per Day]
\textbf{Mathematical Formulation:}
\begin{equation}
\sum_{k \in K} (sr_{ijk} + so_{ijk}) \leq 1 \quad \forall i \in I, \; j \in J
\end{equation}

\textbf{Interpretation:} Each nurse can work at most one shift per day (cannot work both Day and Night shift on the same day).
\end{tcolorbox}

\begin{lstlisting}[language=Python, caption=Constraint 1 Implementation (model.py:L668)]
    for i in I_nurses:
        for j in J_days:
            prob += (
                pulp.lpSum(sr[i][j][k] + so[i][j][k] for k in K_shifts) <= 1,
                f"OneShiftPerDay_{i}_{j}"
            )
\end{lstlisting}

\subsubsection{Constraints 2-5: Min/Max Shifts Per Shift Type}

\begin{tcolorbox}[title=Constraints 2-5: Shift Type Quotas]
\textbf{Mathematical Formulation:}
\begin{align}
\sum_{j \in J} (sr_{ijk} + so_{ijk}) &\geq e_{min}^i \quad \forall i \in I \quad \text{(Min Early)} \\
\sum_{j \in J} (sr_{ijk} + so_{ijk}) &\leq e_{max}^i \quad \forall i \in I \quad \text{(Max Early)} \\
&\text{Similar for Day (D), Late (L), Night (N) shifts}
\end{align}

\textbf{Interpretation:} Each nurse must work within specified min/max bounds for each shift type.
\end{tcolorbox}

\begin{lstlisting}[language=Python, caption=Constraints 2-5 Implementation]
for shift_type, quotas in shift_quotas.items():
    if shift_type in K_shifts:
        shift_min = quotas.get('min', 0)
        shift_max = quotas.get('max', n1)
        
        for i in I_nurses:
            # Minimum shifts of this type
            if shift_min > 0:
                prob += (
                    pulp.lpSum(sr[i][j][shift_type] + so[i][j][shift_type] 
                              for j in J_days) >= shift_min,
                    f"MinShifts_{shift_type}_{i}"
                )
            
            # Maximum shifts of this type
            prob += (
                pulp.lpSum(sr[i][j][shift_type] + so[i][j][shift_type] 
                          for j in J_days) <= shift_max,
                f"MaxShifts_{shift_type}_{i}"
            )
\end{lstlisting}

\subsubsection{Constraint 6: Maximum Total Shifts}

\begin{tcolorbox}[title=Constraint 6: Maximum Total Shifts Per Nurse]
\textbf{Mathematical Formulation:}
\begin{equation}
\sum_{j \in J} \sum_{k \in K} (sr_{ijk} + so_{ijk}) \leq n_1 \quad \forall i \in I
\end{equation}

\textbf{Interpretation:} Each nurse works at most $n_1$ shifts during the planning period (e.g., max 15 shifts in 14 days).
\end{tcolorbox}

\begin{lstlisting}[language=Python, caption=Constraint 6 Implementation (model.py:L771)]
    for i in I_nurses:
        prob += (
            pulp.lpSum(sr[i][j][k] + so[i][j][k] for j in J_days for k in K_shifts) <= n1,
            f"MaxTotalShifts_{i}"
        )
\end{lstlisting}

\subsubsection{Constraint 7: Maximum Night Shifts}

\begin{tcolorbox}[title=Constraint 7: Maximum Night Shifts Per Nurse]
\textbf{Mathematical Formulation:}
\begin{equation}
\sum_{j \in J} (sr_{ijN} + so_{ijN}) \leq n_2 \quad \forall i \in I
\end{equation}

\textbf{Interpretation:} Each nurse works at most $n_2$ night shifts (e.g., max 5 night shifts in 14 days).
\end{tcolorbox}

\begin{lstlisting}[language=Python, caption=Constraint 7 Implementation (model.py:L784)]
    if 'N' in K_shifts:
        for i in I_nurses:
            prob += (
                pulp.lpSum(sr[i][j]['N'] + so[i][j]['N'] for j in J_days) <= n2,
                f"MaxNightShifts_{i}"
            )
\end{lstlisting}

\subsubsection{Constraint 7.5: Weekly Overtime Cap}

\begin{tcolorbox}[title=Constraint 7.5: Weekly Overtime Cap (Case Study Rule)]
\textbf{Paper Reference:} This rule is derived from the case study requirements (Section 5.1) where hospital regulations limit overtime.

\textbf{Mathematical Formulation:}
\begin{equation}
\sum_{j \in Week_w} \sum_{k \in K} so_{ijk} \leq 1 \quad \forall i \in I, w \in W
\end{equation}

\textbf{Interpretation:} Each nurse can work at most \textbf{one} overtime shift per week. This prevents excessive work hours and ensures fair overtime distribution.
\end{tcolorbox}

\begin{lstlisting}[language=Python, caption=Constraint 7.5 Implementation (model.py:L800)]
    if 'MaxOvertimePerWeek' in model_params or True: # Enabled by default
        num_weeks = num_days // 7
        for i in I_nurses:
            for week in range(num_weeks):
                start_day = week * 7 + 1
                end_day = min(num_days, (week + 1) * 7)
                week_days = [d for d in J_days if start_day <= d <= end_day]
                
                prob += (
                    pulp.lpSum(so[i][j][k] for j in week_days for k in K_shifts) <= 1,
                    f"MaxOvertimePerWeek_{i}_week{week}"
                )
\end{lstlisting}

\subsubsection{Constraint 8: Strict Regular Quota (Overtime Distinction)}

\begin{tcolorbox}[title=Constraint 8: Strict Regular Quota]
\textbf{Refinement:} To ensure a clear distinction between regular pay and overtime, $n_3$ is implemented as a \textit{strict quota}. This resolves the "Overtime Paradox" where overtime would otherwise never be used due to its higher cost.

\textbf{Mathematical Formulation:}
\begin{align}
\sum_{j \in J} \sum_{k \in K} sr_{ijk} &\geq n_3 \cdot SR_i \quad \forall i \in I \tag{8a: Min Regular} \\
\sum_{j \in J} \sum_{k \in K} sr_{ijk} &\leq n_3 \cdot SR_i \quad \forall i \in I \tag{8b: Regular Cap}
\end{align}

\textbf{Interpretation:} If a nurse works, they must work \textit{exactly} $n_3$ regular shifts. Any work beyond this limit must be assigned as Overtime ($so_{ijk}$).
\end{tcolorbox}

\begin{lstlisting}[language=Python, caption=Constraint 8 Implementation (model.py:L836)]
    for i in I_nurses:
        # NSS Refinement: Use equality to create a strict quota
        # If nurse works (SR[i]=1), they must work EXACTLY n3 regular shifts
        prob += (
            pulp.lpSum(sr[i][j][k] for j in J_days for k in K_shifts) == n3 * SR[i],
            f"StrictRegularQuota_{i}"
        )
\end{lstlisting}

\subsubsection{Constraint 9: Minimum Complete Weekends Off}

\begin{tcolorbox}[title=Constraint 9: Minimum Complete Weekends Off]
\textbf{Mathematical Formulation:}
\begin{equation}
\sum_{w \in W} \text{WeekendOff}_{iw} \geq n_4 \quad \forall i \in I
\end{equation}

where $\text{WeekendOff}_{iw} = 1$ if nurse $i$ has both Saturday and Sunday off in weekend $w$.

\textbf{Interpretation:} Each nurse must have at least $n_4$ complete weekends (Saturday + Sunday) off.
\end{tcolorbox}

\begin{lstlisting}[language=Python, caption=Constraint 9 Implementation (model.py:L856)]
    if n4 > 0 and weekends:
        for i in I_nurses:
            for w_idx, (sat, sun) in enumerate(weekends):
                # Weekend off only if OFF on BOTH Saturday AND Sunday
                prob += (
                    weekend_off[i][w_idx] <= 1 - pulp.lpSum(sr[i][sat][k] + so[i][sat][k] for k in K_shifts),
                    f"WeekendOff_Sat_{i}_{w_idx}"
                )
                prob += (
                    weekend_off[i][w_idx] <= 1 - pulp.lpSum(sr[i][sun][k] + so[i][sun][k] for k in K_shifts),
                    f"WeekendOff_Sun_{i}_{w_idx}"
                )
            
            # Require minimum number of complete weekends off
            prob += (
                pulp.lpSum(weekend_off[i][w] for w in range(len(weekends))) >= n4,
                f"MinCompleteWeekendsOff_{i}"
            )
\end{lstlisting}

\subsubsection{Constraints 10-13: Night Shift Rest Requirements}

\begin{tcolorbox}[title=Constraints 10-13: Night Shift Rest Rules]
\textbf{Mathematical Formulation:}

\textbf{Constraint 10:} No stand-alone night shifts (must be consecutive)
\begin{equation}
\text{If } sr_{ijN} + so_{ijN} = 1 \text{, then } sr_{i,j-1,N} + so_{i,j-1,N} + sr_{i,j+1,N} + so_{i,j+1,N} \geq 1
\end{equation}

\textbf{Constraint 11:} At least $d$ days off after night shift sequence ends

\textbf{Constraints 12-13:} No Early/Day/Late shifts immediately before or after night shifts

\textbf{Interpretation:} Ensures nurses get adequate rest around night shifts for health and safety.
\end{tcolorbox}

\begin{lstlisting}[language=Python, caption=Constraints 10-11 Implementation (excerpt)]
if night_rest_enabled and 'N' in K_shifts:
    for i in I_nurses:
        for idx, j in enumerate(J_days_sorted):
            # Detect sequence start
            prob += (
                night_sequence_start[i][j] >= 
                    working_night_current - working_night_prev,
                f"NightSeqStart_Detect_{i}_{j}"
            )
            
            # If starting sequence, must work minimum consecutive nights
            for offset in range(1, min_consecutive_nights):
                if idx + offset < len(J_days_sorted):
                    j_future = J_days_sorted[idx + offset]
                    prob += (
                        sr[i][j_future]['N'] + so[i][j_future]['N'] >= 
                            night_sequence_start[i][j],
                        f"MinConsecutiveNights_{i}_{j}_offset{offset}"
                    )
            
            # Days off after night sequence ends
            for offset in range(1, days_off_after_nights + 1):
                if idx + offset < len(J_days_sorted):
                    j_future = J_days_sorted[idx + offset]
                    for k in K_shifts:
                        prob += (
                            sr[i][j_future][k] + so[i][j_future][k] <= 
                                1 - night_sequence_end[i][j],
                            f"DaysOffAfterNights_{i}_{j}_day{offset}_shift{k}"
                        )
\end{lstlisting}
\newpage

\subsubsection{Constraint 14: Stand-Alone Shift Penalty (Soft)}

\begin{tcolorbox}[title=Constraint 14: Stand-Alone Shift Penalty]
\textbf{Mathematical Formulation:}
\begin{equation}
\sum_{k \in K} sr_{i,j-1,k} - \sum_{k \in K} sr_{ijk} + \sum_{k \in K} sr_{i,j+1,k} + dev1_{ij} \geq 0 
\end{equation}
$\forall i \in I, \; j \in \{2, \ldots, |J|-1\}$

\textbf{Interpretation:} Penalizes isolated working days (working day $j$ but not $j-1$ or $j+1$). Encourages consecutive working days.
\end{tcolorbox}

\begin{lstlisting}[language=Python, caption=Constraint 14 Implementation (model.py:L993)]
    for i in I_nurses:
        for idx in range(1, len(J_days_list) - 1):
            j_prev = J_days_list[idx - 1]
            j_curr = J_days_list[idx]
            j_next = J_days_list[idx + 1]
            
            prob += (
                pulp.lpSum(sr[i][j_prev][k] for k in K_shifts) -
                pulp.lpSum(sr[i][j_curr][k] for k in K_shifts) +
                pulp.lpSum(sr[i][j_next][k] for k in K_shifts) +
                dev1[i][j_curr] >= 0,
                f"StandAlonePenalty_{i}_{j_curr}"
            )
\end{lstlisting}

\subsubsection{Constraint 15: Unwanted Shift Pattern Penalty (Soft)}

\begin{tcolorbox}[title=Constraint 15: Unwanted Shift Pattern Penalty]
\textbf{Mathematical Formulation:}
\begin{equation}
sr_{ijk_1} + sr_{i,j+1,k_2} - dev2_{ijk_1} \leq 1 
\end{equation}
$\forall i \in I, \; j \in \{1, \ldots, |J|-1\}, \; (k_1, k_2) \in K'$

where $K' = \{(D,E), (L,E), (L,D), (E,N)\}$ are unwanted consecutive shift combinations.

\textbf{Interpretation:} Penalizes shift sequences with insufficient rest (e.g., Late→Early, Day→Early).
\end{tcolorbox}

\begin{lstlisting}[language=Python, caption=Constraint 15 Implementation (model.py:L1024)]
    for i in I_nurses:
        for idx in range(len(J_days_list) - 1):
            j_curr = J_days_list[idx]
            j_next = J_days_list[idx + 1]
            
            for (k1, k2) in unwanted_patterns:
                # Only add constraint if both shifts exist
                if k1 in K_shifts and k2 in K_shifts:
                    prob += (
                        sr[i][j_curr][k1] + sr[i][j_next][k2] - dev2[i][j_curr][k1] <= 1,
                        f"UnwantedPattern_{i}_{j_curr}_{k1}_{k2}"
                    )
\end{lstlisting}

\subsection{NSS Enhancement: Fatigue Management Constraints}

To prevent nurse burnout and improve patient safety, the model incorporates an exponential fatigue accumulation function based on Jaber et al. (2013).

\subsubsection{Constraint F1: Cumulative Work Hours}
\begin{tcolorbox}[title=Constraint F1: Work Hour Accumulation]
\textbf{Mathematical Formulation:}
\begin{equation}
T_{ij} = T_{i,j-1} + \sum_{k \in K} (sr_{ijk} + so_{ijk}) \cdot \text{duration} \quad \forall i \in I, j \in J
\end{equation}
\textbf{Interpretation:} Tracks the total hours worked by a nurse up to day $j$.
\end{tcolorbox}

\subsubsection{Constraints F2-F5: PWL Fatigue Approximation (SOS2)}
\begin{tcolorbox}[title=Fatigue Approximation with Binary indicators]
\textbf{Refinement:} The exponential function $F(T) = 1 - e^{-\lambda T}$ is non-linear. We use binary segment indicators $y_{ijs}$ to enforce the SOS2 property.

\textbf{Mathematical Formulation:}
\begin{align}
T_{ij} &= \sum_{s \in S} w_{ijs} \cdot B_s \tag{F2: Abscissa} \\
F_{ij} &= \sum_{s \in S} w_{ijs} \cdot V_s \tag{F3: Ordinate} \\
\sum_{s \in S} w_{ijs} &= 1 \tag{F4: Convexity} \\
w_{ijs} &\leq y_{ij,s-1} + y_{ij,s} \tag{F5: SOS2 Link}
\end{align}
where $\sum_s y_{ijs} \leq 2$ and $B_s, V_s$ are PWL data points.
\end{tcolorbox}

\subsubsection{Constraint F6: Patient Safety Threshold}
\begin{tcolorbox}[title=Constraint F6: Fatigue Ceiling]
\textbf{Mathematical Formulation:}
\begin{equation}
F_{ij} \leq F_{max} \quad \forall i, j
\end{equation}
\textbf{Interpretation:} Ensures no nurse exceeds the maximum allowed exhaustion level ($F_{max}=0.70$ by default, which represents significantly increased error risk in clinical settings).
\end{tcolorbox}

\begin{lstlisting}[language=Python, caption=SOS2 Implementation (model.py:L1320)]
        for i in I_nurses:
            for j in J_days:
                # Create binary segment indicators for SOS2
                num_segments = len(breakpoints) - 1  # 8 segments
                y_segment = [pulp.LpVariable(f"SOS2_Segment_{i}_{j}_{s}", cat=pulp.LpBinary)
                            for s in range(num_segments)]
                
                # At most 2 consecutive segments can be active
                prob += (pulp.lpSum(y_segment) <= 2, f"SOS2_MaxTwo_{i}_{j}")
                
                # Link weights to indicator variables
                prob += (pwl_lambda[i, j][0] <= y_segment[0])
                for s in range(1, num_segments):
                    prob += (pwl_lambda[i, j][s] <= y_segment[s-1] + y_segment[s])
                prob += (pwl_lambda[i, j][num_segments] <= y_segment[num_segments - 1])
\end{lstlisting}

\newpage

\subsection{Stage 2 Constraints (Recourse Actions)}

\subsubsection{Constraint 16: Demand Fulfillment (KEY CONSTRAINT)}

\begin{tcolorbox}[title=Constraint 16: Demand Fulfillment with Recourse]
\textbf{Mathematical Formulation:}
\begin{equation}
\sum_{i \in I} (sr_{ijk} + so_{ijk}) + \alpha_{jk}^\omega - \beta_{jk}^\omega \geq R_{jk}^\omega
\end{equation}
$\forall j \in J, \; k \in K, \; \omega \in \Omega$

\textbf{Components:}
\begin{itemize}
    \item $\sum_{i \in I} (sr_{ijk} + so_{ijk})$ = Planned staff (Stage 1)
    \item $\alpha_{jk}^\omega$ = Emergency staff added (Stage 2)
    \item $\beta_{jk}^\omega$ = Shifts cancelled (Stage 2)
    \item $R_{jk}^\omega$ = Actual demand in scenario $\omega$
\end{itemize}

\textbf{Interpretation:} This is the \textbf{critical constraint} linking Stage 1 and Stage 2! After demand is realized, the model can add emergency staff or cancel shifts to meet requirements.
\end{tcolorbox}

\begin{lstlisting}[language=Python, caption=Constraint 16 Implementation (model.py:L1062)]
    for w in W_scenarios:
        for j in J_days:
            for k in K_shifts:
                prob += (
                    pulp.lpSum(sr[i][j][k] + so[i][j][k] for i in I_nurses) +
                    alpha[j][k][w] - beta[j][k][w]
                    >= R_demand.get((j, k, w), 0),
                    f"StaffingMet_{j}_{k}_{w}"
                )
\end{lstlisting}

\subsubsection{Constraint 17: Maximum Emergency Staff Bounds}

\begin{tcolorbox}[title=Constraint 17: Emergency Staff Upper Bound]
\textbf{Mathematical Formulation:}
\begin{equation}
\alpha_{jk}^\omega \leq \alpha^{max} \quad \forall j \in J, \; k \in K, \; \omega \in \Omega
\end{equation}

\textbf{Interpretation:} Limits the maximum number of emergency nurses that can be added per shift (operational/budget constraint).
\end{tcolorbox}

\begin{lstlisting}[language=Python, caption=Constraint 17 Implementation (model.py:L1092)]
    if max_emergency < float('inf'):
        for w in W_scenarios:
            for j in J_days:
                for k in K_shifts:
                    prob += (
                        alpha[j][k][w] <= max_emergency,
                        f"MaxEmergencyStaff_{j}_{k}_{w}"
                    )
\end{lstlisting}

\subsubsection{Constraint 18: Maximum Cancellation Bounds}

\begin{tcolorbox}[title=Constraint 18: Cancellation Upper Bound]
\textbf{Mathematical Formulation:}
\begin{equation}
\beta_{jk}^\omega \leq \beta^{max} \quad \forall j \in J, \; k \in K, \; \omega \in \Omega
\end{equation}

\textbf{Interpretation:} Limits the maximum number of shifts that can be cancelled per shift (operational constraint).
\end{tcolorbox}

\begin{lstlisting}[language=Python, caption=Constraint 18 Implementation (model.py:L1101)]
    if max_cancellations < float('inf'):
        for w in W_scenarios:
            for j in J_days:
                for k in K_shifts:
                    prob += (
                        beta[j][k][w] <= max_cancellations,
                        f"MaxCancellations_{j}_{k}_{w}"
                    )
\end{lstlisting}

\newpage

\subsection{CVaR Risk Management Constraints}

\subsubsection{Constraint 19: CVaR Upper Bound}

\begin{tcolorbox}[title=Constraint 19: CVaR Risk Control]
\textbf{Mathematical Formulation (Rockafellar \& Uryasev, 2000):}
\begin{equation}
\xi + \frac{1}{1-\sigma} \sum_{\omega \in \Omega} p^\omega z^\omega \leq \mu
\end{equation}

\textbf{Components:}
\begin{itemize}
    \item $\xi$ = Value-at-Risk (VaR) threshold
    \item $\sigma$ = Confidence level (e.g., 0.95 for 95\%)
    \item $z^\omega$ = Excess loss in scenario $\omega$ beyond VaR
    \item $\mu$ = User-specified upper bound on CVaR
\end{itemize}

\textbf{Interpretation:} In the worst $(1-\sigma)\%$ of scenarios, the expected shortage must not exceed $\mu$ shifts.

\textbf{Example:} If $\sigma=0.95$ and $\mu=5$: "In the worst 5\% of scenarios, expected shortage $\leq$ 5 shifts"
\end{tcolorbox}

\begin{lstlisting}[language=Python, caption=Constraint 19 Implementation (model.py:L1150)]
    if model_type == "SDM-CVaR":
        prob += (
            xi + (1.0 / (1.0 - sigma)) * pulp.lpSum(
                scenario_probability[w] * z[w] for w in W_scenarios
            ) <= mu,
            "CVaR_Constraint"
        )
\end{lstlisting}

\subsubsection{Constraint 20: Excess Loss Non-Negativity}

\begin{tcolorbox}[title=Constraint 20: Excess Loss Bounds]
\textbf{Mathematical Formulation:}
\begin{equation}
z^\omega \geq 0 \quad \forall \omega \in \Omega
\end{equation}

\textbf{Interpretation:} Automatically enforced by variable bounds. No explicit constraint needed.
\end{tcolorbox}

\begin{lstlisting}[language=Python, caption=Constraint 20 Implementation]
# Automatically enforced by variable definition with lowBound=0
z = pulp.LpVariable.dicts("ExcessLoss_z", 
                         W_scenarios, 
                         lowBound=0,  # This enforces z >= 0
                         cat=pulp.LpContinuous)
\end{lstlisting}

\subsubsection{Constraint 22: Excess Loss Definition}

\begin{tcolorbox}[title=Constraint 22: Excess Loss Calculation]
\textbf{Mathematical Formulation:}
\begin{equation}
z^\omega \geq L^\omega - \xi \quad \forall \omega \in \Omega
\end{equation}

where the loss function is:
\begin{equation}
L^\omega = \sum_{j \in J} \sum_{k \in K} \alpha_{jk}^\omega
\end{equation}

\textbf{Combined with $z^\omega \geq 0$, this gives:}
\begin{equation}
z^\omega = \max(0, L^\omega - \xi)
\end{equation}

\textbf{Interpretation:} 
\begin{itemize}
    \item If $L^\omega \leq \xi$: scenario is not in the tail, $z^\omega = 0$
    \item If $L^\omega > \xi$: scenario is in the worst tail, $z^\omega = L^\omega - \xi$
\end{itemize}
\end{tcolorbox}

\begin{lstlisting}[language=Python, caption=Constraint 22 Implementation (model.py:L1181)]
        for w in W_scenarios:
            # Loss function: total understaffing in scenario w
            loss_function = pulp.lpSum(alpha[j][k][w] 
                                      for j in J_days for k in K_shifts)
            
            prob += (
                z[w] >= loss_function - xi,
                f"ExcessLoss_{w}"
            )
\end{lstlisting}

\newpage

\section{Solution Algorithm}

\subsection{Overall Solution Framework}

\begin{algorithm}[H]
\caption{Two-Stage Stochastic Programming with CVaR}
\begin{algorithmic}[1]
\Require Nurse list $I$, planning horizon $J$, shift types $K$, scenarios $\Omega$
\Ensure Optimal schedule with CVaR risk control
\State \textbf{Input:} Read nurse data, scenario data, model parameters
\State \textbf{Build Stage 1 Variables:} $sr_{ijk}, so_{ijk}, SR_i, SO_i, dev1_{ij}, dev2_{ijk}$
\State \textbf{Build Stage 2 Variables:} $\alpha_{jk}^\omega, \beta_{jk}^\omega$
\State \textbf{Build CVaR Variables:} $\xi, z^\omega$ (if SDM-CVaR)
\State \textbf{Define Objective:} Minimize total expected cost (Equation 1)
\State \textbf{Add Stage 1 Constraints:} Constraints 1-15
\State \textbf{Add Stage 2 Constraints:} Constraints 16-18
\If{model\_type == "SDM-CVaR"}
    \State \textbf{Add CVaR Constraints:} Constraints 19, 20, 22
\EndIf
\State \textbf{Solve:} Use HiGHS/CBC MIP solver
\State \textbf{Extract Solution:} 
\State \quad - Roster: $sr_{ijk}^*, so_{ijk}^*$
\State \quad - Recourse: $\alpha_{jk}^{\omega*}, \beta_{jk}^{\omega*}$
\State \quad - Risk metrics: $\xi^*, z^{\omega*}$
\State \textbf{Return:} Optimal schedule, cost breakdown, risk metrics
\end{algorithmic}
\end{algorithm}

\subsection{Scenario Generation Algorithm}

\begin{algorithm}[H]
\caption{Scenario Generation from Historical Data}
\begin{algorithmic}[1]
\Require Historical demand data, number of scenarios $|\Omega|$
\Ensure Demand scenarios $\{R_{jk}^\omega : \omega \in \Omega\}$
\State \textbf{Collect:} Historical weekly demand patterns over 52 weeks
\State \textbf{Create Pool:} Store all weekly patterns in a pool
\For{$\omega = 1$ to $|\Omega|$}
    \State \textbf{Sample:} Randomly select 4 weekly patterns from pool
    \State \textbf{Concatenate:} Create one monthly scenario (28 days)
    \State \textbf{Convert:} Apply nurse-to-patient ratio to get $R_{jk}^\omega$
    \State \textbf{Set Probability:} $p^\omega = 1/|\Omega|$ (equal probability)
\EndFor
\State \textbf{Return:} Scenario set with equal probabilities
\end{algorithmic}
\end{algorithm}

\newpage

\section{Complete Model Summary}

\subsection{Model Classification}

\begin{itemize}
    \item \textbf{Type:} Two-Stage Stochastic Integer Program
    \item \textbf{Stage 1:} Integer (binary variables for shifts)
    \item \textbf{Stage 2:} Integer recourse (integer add/cancel decisions)
    \item \textbf{Risk Measure:} Conditional Value-at-Risk (CVaR)
    \item \textbf{Solver:} HiGHS (preferred) or CBC (fallback)
\end{itemize}

\subsection{Constraint Count}

\begin{table}[h]
\centering
\begin{tabular}{|l|c|l|}
\hline
\textbf{Constraint} & \textbf{Count} & \textbf{Type} \\
\hline
1. One shift per day & $|I| \times |J|$ & Hard \\
2-5. Shift type quotas & $|I| \times |K| \times 2$ & Hard (optional) \\
6. Max total shifts & $|I|$ & Hard \\
7. Max night shifts & $|I|$ & Hard \\
8. Min regular shifts & $|I|$ & Hard \\
9. Min weekends off & $|I|$ & Hard (optional) \\
10-13. Night rest & $O(|I| \times |J|)$ & Hard (optional) \\
14. Stand-alone penalty & $|I| \times (|J|-2)$ & Soft \\
15. Unwanted patterns & $|I| \times (|J|-1) \times |K'|$ & Soft \\
16. Demand fulfillment & $|J| \times |K| \times |\Omega|$ & Hard \\
17. Max emergency staff & $|J| \times |K| \times |\Omega|$ & Hard (optional) \\
18. Max cancellations & $|J| \times |K| \times |\Omega|$ & Hard (optional) \\
19. CVaR bound & 1 & Hard (SDM-CVaR only) \\
22. Excess loss & $|\Omega|$ & Hard (SDM-CVaR only) \\
\hline
\end{tabular}
\caption{Constraint counts for problem size estimation}
\end{table}

\subsection{Variable Count}

\begin{table}[h]
\centering
\begin{tabular}{|l|c|l|}
\hline
\textbf{Variable} & \textbf{Count} & \textbf{Domain} \\
\hline
$sr_{ijk}$ & $|I| \times |J| \times |K|$ & Binary \\
$so_{ijk}$ & $|I| \times |J| \times |K|$ & Binary \\
$SR_i, SO_i$ & $2|I|$ & Binary \\
$dev1_{ij}$ & $|I| \times |J|$ & Integer $\geq 0$ \\
$dev2_{ijk}$ & $|I| \times |J| \times |K|$ & Integer $\geq 0$ \\
$\alpha_{jk}^\omega$ & $|J| \times |K| \times |\Omega|$ & Integer $\geq 0$ \\
$\beta_{jk}^\omega$ & $|J| \times |K| \times |\Omega|$ & Integer $\geq 0$ \\
$\xi$ & 1 & Continuous \\
$z^\omega$ & $|\Omega|$ & Continuous $\geq 0$ \\
\hline
\end{tabular}
\caption{Variable counts and domains}
\end{table}

\subsection{Complexity Analysis}

For a typical instance:
\begin{itemize}
    \item Nurses: $|I| = 50$
    \item Days: $|J| = 30$
    \item Shifts: $|K| = 4$
    \item Scenarios: $|\Omega| = 20$
\end{itemize}

\textbf{Variable count:}
\begin{align*}
\text{Binary variables (Shifts + indicators)} &\approx 50 \times 30 \times 4 \times 2 + 100 = 12,100 \\
\text{Binary variables (Fatigue SOS2 segments)} &\approx 50 \times 30 \times 8 = 12,000 \\
\text{Total Binary variables} &\approx 24,100 \\
\text{Integer variables (Recourse + Deviations)} &\approx (30 \times 4 \times 20 \times 2) + (50 \times 30 \times 5) = 12,300 \\
\text{Continuous variables} &= 21
\end{align*}

\textbf{Constraint count:} $O(25,000)$ constraints (increased due to SOS2 linkage)

\textbf{Solve time:} 20-30 seconds with HiGHS solver

\section{Implementation Details}

\subsection{Python Package Requirements}

\begin{lstlisting}[language=bash, caption=Required Python Packages]
# Optimization
pulp>=2.7.0
highspy>=1.5.0  # Preferred solver (3-5x faster than CBC)

# Data manipulation
pandas>=2.0.0
numpy>=1.24.0

# Visualization
plotly>=5.17.0
matplotlib>=3.7.0

# Web interface
streamlit>=1.28.0

# Export capabilities
openpyxl>=3.0.0  # Excel export
kaleido>=0.2.1   # PNG export
reportlab>=4.0.0 # PDF export
\end{lstlisting}

\subsection{Solver Selection}

\begin{lstlisting}[language=Python, caption=Automatic Solver Selection (solver\_config.py:L118)]
def create_solver(solver_name, time_limit, mip_gap, verbose=False):
    """
    Create and configure a solver instance.
    Priority: Gurobi > HiGHS > CBC
    """
    if solver_name == 'AUTO' or solver_name is None:
        solver_name = auto_select_solver()
    
    if solver_name == 'GUROBI':
        return pulp.GUROBI(msg=verbose, timeLimit=time_limit, gapRel=mip_gap)
    elif solver_name == 'HiGHS':
        return pulp.HiGHS(msg=verbose, timeLimit=time_limit, gapRel=mip_gap)
    elif solver_name == 'CBC':
        return pulp.PULP_CBC_CMD(msg=verbose, timeLimit=time_limit, gapRel=mip_gap)
\end{lstlisting}

\newpage

\section{Results Interpretation}

\subsection{Output Components}

The model produces the following outputs:

\begin{enumerate}
    \item \textbf{Roster DataFrame:} 
    \begin{itemize}
        \item Rows: Nurses
        \item Columns: Days (D1, D2, ..., D30)
        \item Values: Shift types (E, D, L, N) or "OFF"
    \end{itemize}
    
    \item \textbf{Cost Breakdown:}
    \begin{itemize}
        \item Regular shift costs: $c_1 \sum sr_{ijk}^*$
        \item Overtime costs: $c_2 \sum so_{ijk}^*$
        \item Expected recourse costs: $\sum p^\omega (q^+ \sum \alpha_{jk}^{\omega*} + q^- \sum \beta_{jk}^{\omega*})$
        \item Total cost
    \end{itemize}
    
    \item \textbf{Coverage Analysis:}
    \begin{itemize}
        \item Assigned nurses per day per shift
        \item Demand requirements
        \item Coverage gaps (if any)
    \end{itemize}
    
    \item \textbf{Risk Metrics (SDM-CVaR):}
    \begin{itemize}
        \item VaR value: $\xi^*$
        \item CVaR value: $\xi^* + \frac{1}{1-\sigma} \sum p^\omega z^{\omega*}$
        \item Worst-case shortage scenarios
        \item Shortage distribution across scenarios
    \end{itemize}
    
    \item \textbf{Scenario Analysis:}
    \begin{itemize}
        \item Per-scenario shortages: $\sum_{jk} \alpha_{jk}^{\omega*}$
        \item Per-scenario overages: $\sum_{jk} \beta_{jk}^{\omega*}$
        \item Recourse costs by scenario
    \end{itemize}
\end{enumerate}

\subsection{Quality Metrics}

\begin{equation}
\text{Quality Factor} = 1 - \frac{\sum_t |s_t - d_t|}{2 \sum_t d_t}
\end{equation}

where:
\begin{itemize}
    \item $s_t$ = Number of nurses assigned at time $t$
    \item $d_t$ = Demand (required nurses) at time $t$
\end{itemize}

Quality factor ranges from 0 (worst) to 1 (perfect match).

\newpage

\section{Appendix: Complete Model in Mathematical Notation}

\begin{tcolorbox}[colback=blue!5!white, colframe=blue!75!black, title=I. Objective Function]
\textbf{Minimize:}
\begin{equation*}
\begin{split}
Z = c_1 \sum_{i,j,k} sr_{ijk} &+ c_2 \sum_{i,j,k} so_{ijk} + c_3 \sum_{i,j} dev1_{ij} + c_4 \sum_{i,j,k} dev2_{ijk} \\
& + c_{safety} \sum_{i,j} F_{ij} + \sum_\omega p^\omega \left( q^+ \sum_{j,k} \alpha_{jk}^\omega + q^- \sum_{j,k} \beta_{jk}^\omega \right)
\end{split}
\end{equation*}
\end{tcolorbox}

\begin{tcolorbox}[colback=gray!5!white, colframe=gray!75!black, title=II. Stage 1: Base Staffing \& Patterns]
\begin{align*}
\sum_k (sr_{ijk} + so_{ijk}) &\leq 1 && \forall i,j \tag{C1} \\
e_{min} \leq \sum_j (sr_{ijE} + so_{ijE}) &\leq e_{max} && \forall i \tag{C2-3} \\
d_{min} \leq \sum_j (sr_{ijD} + so_{ijD}) &\leq d_{max} && \forall i \tag{C4-5} \\
\sum_{j,k} (sr_{ijk} + so_{ijk}) &\leq n_1 && \forall i \tag{C6} \\
\sum_j (sr_{ijN} + so_{ijN}) &\leq n_2 && \forall i \tag{C7} \\
\sum_{j \in Week_w} \sum_k so_{ijk} &\leq 1 && \forall i, w \tag{C7.5} \\
\sum_{j,k} sr_{ijk} &= n_3 \cdot SR_i && \forall i \tag{C8} \\
\sum_w \text{WeekendOff}_{iw} &\geq n_4 && \forall i \tag{C9}
\end{align*}
\end{tcolorbox}

\begin{tcolorbox}[colback=green!5!white, colframe=green!75!black, title=III. Stage 1: Advanced Sequencing \& Penalties]
\begin{align*}
\text{NightSeqStart}_{ij} \geq (sr_{ijN}&+so_{ijN}) - (sr_{i,j-1,N}+so_{i,j-1,N}) && \forall i,j \tag{C10a} \\
sr_{i,j+m,N} + so_{i,j+m,N} &\geq \text{NightSeqStart}_{ij} && m \in \{1, \dots, min\_consec\} \tag{C10b} \\
\sum_k (sr_{i,j+m,k} + so_{i,j+m,k}) &\leq 1 - \text{NightSeqEnd}_{ij} && m \in \{1, \dots, days\_off\} \tag{C11} \\
\sum_k sr_{i,j-1,k} - \sum_k sr_{ijk} &+ \sum_k sr_{i,j+1,k} + dev1_{ij} \geq 0 && \forall i,j \tag{C14} \\
sr_{ijk_1} + sr_{i,j+1,k_2} &- dev2_{ijk_1} \leq 1 && \forall i,j, (k_1,k_2) \in K' \tag{C15}
\end{align*}
\end{tcolorbox}

\begin{tcolorbox}[colback=orange!5!white, colframe=orange!75!black, title=IV. Stage 1: Fatigue Management (NSS Contribution)]
\begin{align*}
T_{ij} = T_{i,j-1} + &\sum_k (sr_{ijk} + so_{ijk}) \cdot \text{duration} && \forall i,j \tag{F1} \\
\sum_s w_{ijs} &= 1 && \forall i,j \tag{F2} \\
T_{ij} = \sum_s w_{ijs}& B_s, \quad F_{ij} = \sum_s w_{ijs} V_s && \forall i,j \tag{F3-4} \\
w_{ijs} \leq y_{ij,s-1} &+ y_{ij,s}, \quad \sum_s y_{ijs} \leq 2 && \forall i,j,s \tag{F5} \\
F_{ij} &\leq F_{max} && \forall i,j \tag{F6}
\end{align*}
\end{tcolorbox}

\begin{tcolorbox}[colback=red!5!white, colframe=red!75!black, title=V. Stage 2 \& Risk Management (CVaR)]
\begin{align*}
\sum_i (sr_{ijk} + so_{ijk}) + \alpha_{jk}^\omega &- \beta_{jk}^\omega \geq R_{jk}^\omega && \forall j,k,\omega \tag{C16} \\
\alpha_{jk}^\omega &\leq \alpha^{max}, \quad \beta_{jk}^\omega \leq \beta^{max} && \forall j,k,\omega \tag{C17-18} \\
\xi + \frac{1}{1-\sigma} \sum_\omega p^\omega z^\omega &\leq \mu && \text{(SDM-CVaR)} \tag{C19} \\
z^\omega \geq \sum_{j,k} \alpha_{jk}^\omega &- \xi, \quad z^\omega \geq 0 && \forall \omega \tag{C20,22}
\end{align*}
\vspace{5pt}
\textbf{Variable Domains:}
\begin{align*}
sr_{ijk}, so_{ijk}, SR_i, SO_i, y_{ijs} &\in \{0,1\}, \quad \alpha_{jk}^\omega, \beta_{jk}^\omega, dev1_{ij}, dev2_{ijk} \in \mathbb{Z}_+ \\
w_{ijs}, F_{ij}, T_{ij}, \xi &\in \mathbb{R}
\end{align*}
\end{tcolorbox}

\newpage

\section{Complete Paper-to-Code Mapping}

\subsection{Equation-by-Equation Comparison: Objective Function}

\begin{longtable}{|p{4cm}|p{5cm}|p{5cm}|}
\caption{Objective Function Components Mapping} \\
\hline
\textbf{Paper Component} & \textbf{Mathematical Form} & \textbf{Code Location} \\
\hline
\endfirsthead

\multicolumn{3}{c}{\tablename\ \thetable\ -- Continued from previous page} \\
\hline
\textbf{Paper Component} & \textbf{Mathematical Form} & \textbf{Code Location} \\
\hline
\endhead

\hline
\multicolumn{3}{r}{Continued on next page} \\
\endfoot

\hline
\endlastfoot

Regular shift cost & $c_1 \sum_{ijk} sr_{ijk}$ & \texttt{model.py:617} \\
\hline
Overtime shift cost & $c_2 \sum_{ijk} so_{ijk}$ & \texttt{model.py:619} \\
\hline
Standalone penalty & $c_3 \sum_{ij} dev1_{ij}$ & \texttt{model.py:624} \\
\hline
Unwanted pattern penalty & $c_4 \sum_{ijk} dev2_{ijk}$ & \texttt{model.py:625} \\
\hline
Fatigue Safety Penalty & $c_{safety} \sum F_{ij}$ & \texttt{model.py:639} \\
\hline
Emergency staff cost & $q^+ \sum_{jk\omega} \alpha_{jk}^\omega$ & \texttt{model.py:629-636} \\
\hline
Cancellation cost & $q^- \sum_{jk\omega} \beta_{jk}^\omega$ & \texttt{model.py:629-636} \\
\hline
Expected recourse & $\sum_\omega p^\omega (\dots)$ & \texttt{model.py:629-636} \\
\hline
\end{longtable}

\subsection{Equation-by-Equation Comparison: Stage 1 Constraints}

\begin{longtable}{|p{4cm}|p{5cm}|p{5cm}|}
\caption{Stage 1 Constraints Mapping} \\
\hline
\textbf{Paper Component} & \textbf{Mathematical Form} & \textbf{Code Location} \\
\hline
\endfirsthead

\multicolumn{3}{c}{\tablename\ \thetable\ -- Continued from previous page} \\
\hline
\textbf{Paper Component} & \textbf{Mathematical Form} & \textbf{Code Location} \\
\hline
\endhead

\hline
\multicolumn{3}{r}{Continued on next page} \\
\endfoot

\hline
\endlastfoot

C1: One shift/day & $\sum_k (sr+so) \leq 1$ \newline $\forall i \in I, j \in J$ & \texttt{model.py:668} \\
\hline
C2-5: Shift quotas & $n_{min} \leq \sum_j sr_{ijk} \leq n_{max}$ \newline $\forall i \in I, k \in K$ & \texttt{model.py:678} \\
\hline
C6: Max total shifts & $\sum_{jk} (sr+so) \leq n_1$ \newline $\forall i \in I$ & \texttt{model.py:771} \\
\hline
C7: Max night shifts & $\sum_j sr_{ijN} \leq n_2$ \newline $\forall i \in I$ & \texttt{model.py:784} \\
\hline
C7.5: Weekly OT Cap & $\sum_{j \in Week} so \leq 1$ & \texttt{model.py:800} \\
\hline
C8: Min regular shifts & $\sum_{jk} sr_{ijk} \geq n_3$ \newline $\forall i \in I$ & \texttt{model.py:838} \\
\hline
C8b: Regular Quota Cap & $\sum_{jk} sr_{ijk} \leq n_3$ \newline $\forall i \in I$ & \texttt{model.py:840} \\
\hline
F1: Weariness Tracking & $T_{ij} = T_{i,j-1} + \dots$ & \texttt{model.py:1225} \\
\hline
F2-F5: PWL Fatigue & $F_{ij} = \sum w_{ijs} V_s$ & \texttt{model.py:1253} \\
\hline
F6: Safety Threshold & $F_{ij} \leq F_{max}$ & \texttt{model.py:1377} \\
\hline
C9: Min weekends off & $\sum_w W_{iw} \geq n_4$ \newline $\forall i \in I$ & \texttt{model.py:856} \\
\hline
C10: No standalone nights & $sr_{ijN}=1 \Rightarrow$ \newline $sr_{i,j-1,N}+sr_{i,j+1,N}\geq 1$ & \texttt{model.py:917} \\
\hline
C11: Days off after nights & Consecutive nights $\Rightarrow$ \newline $d$ days off after & \texttt{model.py:942} \\
\hline
C14: Standalone penalty & $\sum_k sr_{j-1} - \sum_k sr_j$ \newline $+ \sum_k sr_{j+1} + dev1 \geq 0$ & \texttt{model.py:993} \\
\hline
C15: Unwanted patterns & $sr_{ijk_1} + sr_{i,j+1,k_2}$ \newline $- dev2 \leq 1$ & \texttt{model.py:1024} \\
\hline
\end{longtable}

\subsection{Equation-by-Equation Comparison: Stage 2 and CVaR}

\begin{longtable}{|p{4cm}|p{5cm}|p{5cm}|}
\caption{Stage 2 and CVaR Constraints Mapping} \\
\hline
\textbf{Paper Component} & \textbf{Mathematical Form} & \textbf{Code Location} \\
\hline
\endfirsthead

\multicolumn{3}{c}{\tablename\ \thetable\ -- Continued from previous page} \\
\hline
\textbf{Paper Component} & \textbf{Mathematical Form} & \textbf{Code Location} \\
\hline
\endhead

\hline
\multicolumn{3}{r}{Continued on next page} \\
\endfoot

\hline
\endlastfoot

C16: Demand fulfillment & $\sum_i (sr+so) + \alpha - \beta$ \newline $\geq R_{jk}^\omega$ & \texttt{model.py:1062} \\
\hline
C17: Max emergency & $\alpha_{jk}^\omega \leq \alpha^{max}$ \newline $\forall j,k,\omega$ & \texttt{model.py:1092} \\
\hline
C18: Max cancellations & $\beta_{jk}^\omega \leq \beta^{max}$ \newline $\forall j,k,\omega$ & \texttt{model.py:1101} \\
\hline
C19: CVaR bound & $\xi + \frac{1}{1-\sigma} \sum p^\omega z^\omega$ \newline $\leq \mu$ & \texttt{model.py:1150} \\
\hline
C20: Excess loss $\geq 0$ & $z^\omega \geq 0$ \newline $\forall \omega$ & \texttt{model.py:283} \\
\hline
C22: Excess loss def & $z^\omega \geq \sum_{jk} \alpha^\omega - \xi$ \newline $\forall \omega$ & \texttt{model.py:1181} \\
\hline
\end{longtable}

\subsection{Variable Definitions Comparison}

\begin{longtable}{|l|l|l|l|}
\caption{Variable Naming Correspondence} \\
\hline
\textbf{Variable} & \textbf{Paper Symbol} & \textbf{Code Name} & \textbf{Code Line} \\
\hline
\endfirsthead

\multicolumn{4}{c}{\tablename\ \thetable\ -- Continued from previous page} \\
\hline
\textbf{Variable} & \textbf{Paper Symbol} & \textbf{Code Name} & \textbf{Code Line} \\
\hline
\endhead

\hline
\multicolumn{4}{r}{Continued on next page} \\
\footnotetext{Note: $sr$ and $so$ indicate regular and overtime shifts respectively.}
\endfoot

\hline
\endlastfoot

\multicolumn{4}{|c|}{\textbf{STAGE 1 VARIABLES}} \\
\hline
Regular shift & $sr_{ijk} \in \{0,1\}$ & \texttt{sr[i][j][k]} & \texttt{model.py:403} \\
\hline
Overtime shift & $so_{ijk} \in \{0,1\}$ & \texttt{so[i][j][k]} & \texttt{model.py:409} \\
\hline
Regular indicator & $SR_i \in \{0,1\}$ & \texttt{SR[i]} & \texttt{model.py:415} \\
\hline
Overtime indicator & $SO_i \in \{0,1\}$ & \texttt{SO[i]} & \texttt{model.py:416} \\
\hline
Cumulative Fatigue & $F_{ij} \in [0, F_{max}]$ & \texttt{F[i][j]} & \texttt{model.py:521} \\
\hline
Cumulative Hours & $T_{ij} \in \mathbb{R}_+$ & \texttt{T[i][j]} & \texttt{model.py:529} \\
\hline
PWL Weights & $w_{ijs} \in [0,1]$ & \texttt{pwl\_lambda[i,j]} & \texttt{model.py:506} \\
\hline
SOS2 Indicators & $y_{ijs} \in \{0,1\}$ & \texttt{y\_segment} & \texttt{model.py:1324} \\
\hline
Standalone deviation & $dev1_{ij} \in \mathbb{Z}_+$ & \texttt{dev1[i][j]} & \texttt{model.py:427} \\
\hline
Pattern deviation & $dev2_{ijk} \in \mathbb{Z}_+$ & \texttt{dev2[i][j][k]} & \texttt{model.py:435} \\
\hline
Weekend off & $W_{iw} \in \{0,1\}$ & \texttt{weekend\_off[i][w]} & \texttt{model.py:478} \\
\hline
\multicolumn{4}{|c|}{\textbf{STAGE 2 VARIABLES}} \\
\hline
Shortage & $\alpha_{jk}^\omega \in \mathbb{Z}_+$ & \texttt{alpha[j][k][w]} & \texttt{model.py:555} \\
\hline
Surplus & $\beta_{jk}^\omega \in \mathbb{Z}_+$ & \texttt{beta[j][k][w]} & \texttt{model.py:562} \\
\hline
\multicolumn{4}{|c|}{\textbf{CVAR VARIABLES}} \\
\hline
Value-at-Risk & $\xi \in \mathbb{R}$ & \texttt{xi} & \texttt{model.py:581} \\
\hline
Excess loss & $z^\omega \in \mathbb{R}_+$ & \texttt{z[w]} & \texttt{model.py:586} \\
\hline
\end{longtable}

\end{document}